\section{Validation Audits}
\label{app:validation_audits_v18}

The validation design has to address two concerns. First,
frequent participants have more chances to appear near any
defendant, so raw co-bidding rates are not enough. Second,
item-level labels can mechanically reuse the same participation
history that constructs the score. The main text reports the
headline results; this appendix records the audit logic.

\subsection{Participation-Stratified Permutation}
\label{app:permutation_v18}

The conservative CADE benchmark fixes the direct-defendant list
using cases adjudicated before the end of
$\valConservativeCutoff$. Within the always-loser stratum, we
then compare the observed frequent-loser share among cobidders
with a permutation distribution that preserves participation
volume. Specifically, firms are grouped by participation-volume
quartile, cobidder status is randomly permuted within those
quartiles, and the frequent-loser share among placebo cobidders
is recomputed $\valPermPilot$ times.

The observed share is $\valCADEenrich$, compared with a
permutation mean of $\valCADEbase$. The ratio is
$\valMultThreeTwo\!\times$, with $p<0.001$, and the conservative
ROC AUC is $\valAUCprePost$ (95\% CI $\valAUCprePostCI$). The
permutation does not make cobidder status causal or legal proof.
It answers the narrower exposure concern: the frequent-loser
concentration near CADE defendants is larger than participation
volume alone predicts.

\subsection{Leakage Decomposition}
\label{app:leakage_v18}

The in-sample item-level AUC is intentionally not treated as the
main validation number. It is $\valLeakRefAUC$, and part of that
performance is mechanical: the same firms can help define both
the participation score and the item-level cobidder label. The
audit decomposes that headline into a structural component and a
leakage component.

\begin{table}[htbp]
\centering
\caption{Leakage Audit}
\label{tab:leakage_audit_v18}
\footnotesize
\begin{adjustbox}{max width=\textwidth,center}
\begin{threeparttable}
\begin{tabular}{p{4.4cm}p{5.4cm}cc}
\toprule
Audit & What changes & AUC & 95\% CI \\
\midrule
In-sample reference & Full-sample score predicts any-cobidder item label
& $\valLeakRefAUC$ & $\valLeakRefAUCCI$ \\
Direct-defendant scope check & Same score predicts direct-CADE item label
& $\valLeakDirectAUC$ & $\valLeakDirectAUCCI$ \\
Cobidder-firm holdout & Held-out cobidders are removed from score construction within folds
& $\valLeakCVAUC$ & $\valLeakCVAUCCI$ \\
Temporal holdout & $2009$--$2016$ score predicts $2017$--$2019$ cobidder labels
& $\valAUCFLfirmTemp$ & $\valAUCFLfirmTempCI$ \\
Temporal direct-defendant check & Same temporal score predicts direct-CADE labels
& $\valLeakDirectTempAUC$ & $\valLeakDirectTempAUCCI$ \\
\bottomrule
\end{tabular}
\begin{tablenotes}
\footnotesize
\item \textit{Notes:} The reference AUC is partly tautological.
The firm-holdout and temporal rows isolate the component that
survives when labels are withheld from score construction or
placed out of time. Direct-defendant rows check whether the
loser-side statistic accidentally behaves like a general
membership detector.
\end{tablenotes}
\end{threeparttable}
\end{adjustbox}
\end{table}

The surviving component is large: AUC remains in the
$\valLeakStructLow$--$\valLeakStructHigh$ range under the
holdout audits. The lost component, roughly
$\valLeakTautLow$--$\valLeakTautHigh$ AUC, is disclosed as
mechanical. That disclosure is important because the main claim
is not that a record-layer statistic can classify all cartel
conduct perfectly. It is that a cheap record-layer statistic
retains meaningful ranking power after the most obvious
mechanical channels are removed.

\subsection{Direct-Defendant Scope Check}
\label{app:direct_scope_v18}

The direct-defendant result is a falsification of overreach. If
the statistic were a generic cartel-membership detector, it
should classify direct CADE defendants as well as cobidders. It
does not. The direct-defendant AUC is $\valAUCdirectStd$, close
to random classification. The result is not a weakness to be
hidden. It is the expected signature of a loser-side screen: the
screen ranks firms that repeatedly lose, while direct defendants
are legal anchors that often appear on the winner side of the
arrangement.

\subsection{Year-by-Year Holdout}
\label{app:year_holdout_v18}

The rolling-origin exercise trains the score on years preceding
each test year and evaluates cobidder discrimination in the test
year. Table~\ref{tab:temporal_holdout_year_v18} reports the
year-by-year decomposition behind the temporal-holdout figure in
the main text.

\begin{table}[htbp]
\centering
\caption{Temporal-Holdout AUC by Test Year}
\label{tab:temporal_holdout_year_v18}
\footnotesize
\begin{threeparttable}
\begin{tabular}{ccccc}
\toprule
Test year & AUC & 95\% CI & $N$ firms & $N$ CADE \\
\midrule
$2014$ & $0.819$ & $[0.779,\ 0.860]$ & $8{,}631$ & $88$ \\
$2015$ & $0.817$ & $[0.776,\ 0.858]$ & $10{,}122$ & $102$ \\
$2016$ & $0.851$ & $[0.821,\ 0.881]$ & $11{,}699$ & $127$ \\
$2017$ & $0.862$ & $[0.832,\ 0.891]$ & $13{,}057$ & $148$ \\
$2018$ & $0.897$ & $[0.877,\ 0.918]$ & $14{,}391$ & $161$ \\
$2019$ & $0.922$ & $[0.912,\ 0.933]$ & $15{,}597$ & $181$ \\
\bottomrule
\end{tabular}
\begin{tablenotes}
\footnotesize
\item \textit{Notes:} Each row computes the screening score using
participation observed before the test year and evaluates AUC
against the cobidder label in that test year.
\end{tablenotes}
\end{threeparttable}
\end{table}

The year-by-year pattern should not be read as a claim that the
environment becomes progressively easier to police. The samples
and adjudication links change across years. The relevant point is
more limited: the score remains separated from random
classification in each rolling-origin test year.
